\section{About the Committee}

\subsection{Purpose}

The North Atlantic Organization Treaty was founded in 1949 with three main purposes:

\begin{enumerate}
  \item Deterring Soviet expansionism
  \item Forbidding the revival of nationalist militarism in Europe through a strong North American presence on the continent
  \item Encouraging European political integration.
\end{enumerate}

The aftermath of World War II saw much of Europe devastated in a way that is now difficult to envision. Approximately 36.5 million Europeans had died in the conflict, 19 million of them civilians. Shortly after the conflict, Communists, aided by the Soviet Union, were threatening elected governments across Europe as happened in Czechoslovakia in 1948\footnote{North Atlantic Treaty Organization. (n.d.). A short history of NATO. NATO. \url{https://www.nato.int/en/about-us/nato-history/a-short-history-of-nato}}.

The establishment of the Marshall plan, funded by the United States, gave a solid economic base to Western European democracies to come together to implement projects that promoted military cooperation and collective defence. After multiple negotiations, in April, 1949, the treaty was signed, being Article 5 one of the most important ones. The latter establishes that \emph{an armed attack against one or more of them… shall be considered an attack against them all} and that following such an attack, each Ally would take \emph{such action as it deems necessary, including the use of armed force} in response. The only time this Article was invoked was as a response to the 9/11 terrorist attacks in the United States in 2001.\footnote{North Atlantic Treaty Organization. (n.d.). A short history of NATO. NATO. \url{https://www.nato.int/en/about-us/nato-history/a-short-history-of-nato}}

\subsection{Mandate and Treaty}

The purpose of NATO is therefore the one to guarantee the freedom and security of its member states through political and military means.

On one hand, political means refers to the promotion of democratic values and the provision of a structure to consult and cooperate on defence and security. On the other hand, military means refer to the last resort available if diplomatic means fail. This is inherent to NATO’S defensive alliance, since its main purpose is to keep its members safe and secure by maintaining a credible deterrence and defence posture.\footnote{North Atlantic Treaty Organization. (n.d.). What is NATO? NATO. \url{https://www.nato.int/en/what-is-nato}}

It is very important to highlight that \textbf{all decisions at NATO are taken by consensus. This means that in order to pass a resolution, all 32 member states must agree. Delegates must take this into account when negotiating and cooperating among different blocs.}

The most relevant articles of the treaty are summarised below. However, we encourage you to read the full treaty to acquire a deeper understanding of NATO’s capabilities and obligations.

\begin{guidetable}[Y{0.3}Y{1.7}]{2}
  \emph{ARTICLE 1} & \emph{Settle any international dispute in which they may be involved by peaceful means in such a manner that international peace and security and justice are not endangered, and to refrain in their international relations from the threat or use of force in any manner inconsistent with the purposes of the United Nations.} \\
  \midrule
  \emph{ARTICLE 2} & \emph{Contribute toward the further development of peaceful and friendly international relations.} \newline \emph{Seek to eliminate conflict in their international economic policies and will encourage economic collaboration between any or all of them.} \\
  \midrule
  \emph{ARTICLE 3} & \emph{Parties, separately and jointly will maintain and develop their individual and collective capacity to resist armed attack} \\
  \midrule
  \emph{ARTICLE 4} & \emph{Consult the territorial integrity, political independence or security of any of the Parties is threatened.} \\
  \midrule
  \emph{ARTICLE 5} & \emph{The Parties agree that an armed attack against one or more of them in Europe or North America shall be considered an attack against them all If such an armed attack occurs, each of them, in exercise of the right of individual or collective self-defence recognised by Article 51 of the Charter of the United Nations, will assist the Party or Parties so attacked by taking forthwith, individually and in concert with the other Parties, including the use of armed force, to restore and maintain the security of the North Atlantic area.} \\
  \midrule
  \emph{ARTICLE 6} & \emph{For the purpose of Article 5, an armed attack on one or more of the Parties is deemed to include an armed attack: On the territory of any of the Parties in Europe or North America, on the Algerian Departments of France, on} \emph{the territory of Turkey or on the Islands under the jurisdiction of any of the Parties in the North Atlantic area north of the Tropic of Cancer;} \newline \emph{On the forces, vessels, or aircraft of any of the Parties, when in or over these territories or any other area in Europe in which occupation forces of any of the Parties were stationed on the date when the Treaty entered into force or the Mediterranean Sea or the North Atlantic area north of the Tropic of Cancer.} \\
  \midrule
  \emph{ARTICLE 8} & \emph{Each Party declares that none of the international engagements now in force between it and any other of the Parties or any third State is in conflict with the provisions of this Treaty, and undertakes not to enter into any international engagement in conflict with this Treaty.} \\
  \midrule
  \emph{ARTICLE 10} & \emph{The Parties may, by unanimous agreement, invite any other European State to contribute to the security of the North Atlantic area to accede to this Treaty.} \\
  \midrule
  \emph{ARTICLE 12} & \emph{After the Treaty has been in force for ten years, or at any time thereafter, the Parties shall, if any of them so requests, consult together for the purpose of reviewing the Treaty.} \\
\end{guidetable}

\subsection{Partners and non-member states}

Even though the NATO Alliance only applies to the member states, NATO cooperates with multiple partners, primarily based on their geographic region. They contribute to many of NATO's core activities, from shaping policy to building defence capacity, developing interoperability and managing crises.

\textbf{Partnership for Peace (PfP)}: bilateral cooperation between NATO and individual partner countries in the Euro-Atlantic area.

The PfP was established in 1994 to enable participants to develop an individual relationship with NATO, choosing their own priorities for cooperation, and the level and pace of progress. The countries involved are: Armenia, Azerbaijan, Georgia, Kazakhstan, Malta, Serbia, Tajikistan, Ukraine, Bosnia and Herzegovina, Ireland, Kyrgyzstan, Republic of Moldova, Switzerland, Turkmenistan, Uzbekistan.

\textbf{Mediterranean Dialogue:} aims to contribute to security and stability in the wider Mediterranean region.

The MD's core principles are: non-discrimination, self-differentiation, two-way engagement, non-imposition, diversity, and complementarity to other international initiatives in the region. The countries involved are: Algeria, Israel, Mauritania, Tunisia, Egypt, Jordan, Morocco.\footnote{North Atlantic Treaty Organization. (n.d.). Mediterranean Dialogue. NATO. \url{https://www.nato.int/en/what-we-do/partnerships-and-cooperation/mediterranean-dialogue}}

\textbf{Istanbul Cooperation Initiative:} offers non-NATO countries in the Middle East region the opportunity to cooperate with NATO.

The countries that belong to this alliance are: Bahrain, Qatar, Kuwait, United Arab Emirates.

\textbf{Partners across the global:} partnerships outside the formal structures, on an individual basis. Those countries include: Australia, Iraq, Mongolia, Pakistan, Colombia, Japan, New Zealand, Republic of Korea.

\textbf{International organizations:} European Union, United Nations, Organization for Security and Cooperation in Europe.
